\documentclass[11pt]{article}

\usepackage{amsmath,amssymb,amsfonts,bm}
\usepackage[a4paper,margin=1in]{geometry}
\usepackage{mathtools}
\usepackage{hyperref}

% Chebyshev
\newcommand{\T}{\mathcal{T}}
\newcommand{\U}{\mathcal{U}}

\title{Exact Generating Functions and First-Passage Statistics on a 1D Ring\\
with One \emph{Directed} Shortcut (defect technique; nearest-neighbour $k=1$; $K=2$ defect sites)\\
\large Notation: ($\underline{B},\underline{A},\underline{X},\eta$)}
\author{}
\date{}

\begin{document}
\maketitle

\section*{Notation and model}

\paragraph{Ring lattice.}
We consider a discrete-time random walk on a 1D ring (cycle) of size $N$ with periodic boundary conditions (PBC).
Sites are labeled $n\in\{0,1,\dots,N-1\}$ and all indices are understood modulo $N$.

\paragraph{Transition matrices and convention.}
We use the convention
\[
\begin{aligned}
\Pr(X_{t+1}=n\mid X_t=n') &= B_{n,n'} \qquad \text{(defect-free ring)},\\
\Pr(X_{t+1}=n\mid X_t=n') &= A_{n,n'} \qquad \text{(with defect)}.
\end{aligned}
\]
Thus, the master equation is $P(n,t+1)=\sum_{n'} B_{n,n'} P(n',t)$ for the defect-free process.
Both $\underline{B}$ and $\underline{A}$ are column-stochastic: $\sum_n B_{n,n'}=\sum_n A_{n,n'}=1$.

\paragraph{Baseline (defect-free) lazy random walk.}
Let $q\in(0,1)$ be the \emph{total} nearest-neighbour ($k=1$) jump probability. Then
\[
B_{n,n}=1-q,\qquad B_{n\pm 1,n}=\frac{q}{2},\qquad B_{m,n}=0\text{ otherwise}.
\]

\paragraph{Remark on the notation $k$ vs $K$.}
This document considers the nearest-neighbour ring (jumps only to $n\pm 1$, i.e.\ $k=1$).
The symbol $K=2$ is used here \emph{only} to count the two defect sites at the shortcut endpoints $(u,v)$;
it should not be confused with the connectivity parameter $K=2k$ used in other ring-lattice/small-world models.

\paragraph{Single directed shortcut.}
We add \emph{one directed} (unidirectional) shortcut from a site $u$ to a site $v$.
Its probability is denoted $p$ and is \emph{taken from the lazy weight at $u$}. Concretely, only the
outgoing probabilities from $u$ are changed:
\[
A_{u,u}=1-q-p,\qquad A_{u\pm 1,u}=\frac q2,\qquad A_{v,u}=B_{v,u}+p,\qquad
A_{n,u}=B_{n,u}\ \text{otherwise},
\]
and for all $n'\neq u$ we keep $A_{n,n'}=B_{n,n'}$.
The constraint is
\[
0\le p\le 1-q
\]
so that $A_{u,u}\ge 0$.

\paragraph{Defect matrix and defect parameters.}
We encode the modifications by the (sparse) defect matrix
\[
\underline{X}\equiv \underline{B}-\underline{A}.
\]
In the defect-technique formulas of the Supplement, the defect parameters associated with a pair $(u_i,v_i)$ are
written $\eta_{u_i,v_i}$ and $\eta_{v_i,u_i}$ and correspond to the off-diagonal entries of $\underline{X}$:
\[
\eta_{u,v}\equiv X_{v,u}=B_{v,u}-A_{v,u},\qquad \eta_{v,u}\equiv X_{u,v}=B_{u,v}-A_{u,v}.
\]
For a \emph{single directed} shortcut $u\to v$ of probability $p$ (built from the lazy weight at $u$) we have
\[
\boxed{\eta_{u,v}=-p,\qquad \eta_{v,u}=0,}
\]
while the corresponding diagonal change at $u$ is implied by probability conservation:
$X_{u,u}=B_{u,u}-A_{u,u}=p=-\eta_{u,v}$.

\paragraph{Propagators and generating functions.}
\begin{itemize}
\item Defect-free (ring) propagator: $Q_{n_0}(n,t)=\Pr(X_t=n\mid X_0=n_0)$ under $\underline{B}$.
\item Defective (network) propagator: $S_{n_0}(n,t)=\Pr(X_t=n\mid X_0=n_0)$ under $\underline{A}$.
\item $z$-transform: $\tilde f(z)=\sum_{t=0}^\infty f(t)\,z^t$.
\end{itemize}
We denote $\tilde Q_{n_0}(n,z)=\sum_{t\ge 0} Q_{n_0}(n,t)z^t$ and similarly $\tilde S_{n_0}(n,z)$.

\paragraph{Ring distance.}
For any two sites $a,b$ define
\[
d(a,b)\equiv \min\bigl(|a-b|,\;N-|a-b|\bigr)\in\Bigl\{0,1,\dots,\Bigl\lfloor \frac N2\Bigr\rfloor\Bigr\}.
\]

\paragraph{Chebyshev polynomials.}
We use
\[
\T_n(x)=\cos\!\bigl(n\arccos x\bigr),\qquad
\U_n(x)=\frac{\sin\!\bigl((n+1)\arccos x\bigr)}{\sin(\arccos x)}.
\]

\section{Part 1 --- Homogeneous (defect-free) propagator $\tilde Q_{n_0}(n,z)$ on the ring}

\subsection{Generating function on a finite ring}

The defect-free master equation is
\begin{equation}
Q_{n_0}(n,t+1)=\frac q2\,Q_{n_0}(n-1,t)+\frac q2\,Q_{n_0}(n+1,t)+(1-q)\,Q_{n_0}(n,t),
\label{eq:master_ring}
\end{equation}
with $Q_{n_0}(n,0)=\delta_{n,n_0}$ and PBC.

Taking the $z$-transform of \eqref{eq:master_ring} proceeds as follows.
Multiply both sides of \eqref{eq:master_ring} by $z^t$ and sum from $t=0$ to $\infty$:
\begin{align*}
\sum_{t=0}^{\infty} Q_{n_0}(n,t+1) z^t
&=\frac q2\sum_{t=0}^{\infty} Q_{n_0}(n-1,t) z^t
 +\frac q2\sum_{t=0}^{\infty} Q_{n_0}(n+1,t) z^t
 +(1-q)\sum_{t=0}^{\infty} Q_{n_0}(n,t) z^t.
\end{align*}
The three sums on the right are just the generating functions:
\[
\sum_{t=0}^{\infty} Q_{n_0}(n\pm 1,t) z^t=\tilde Q_{n_0}(n\pm 1,z),
\qquad
\sum_{t=0}^{\infty} Q_{n_0}(n,t) z^t=\tilde Q_{n_0}(n,z).
\]
For the time-shifted left-hand side we use the identity
\begin{align*}
\sum_{t=0}^{\infty} Q_{n_0}(n,t+1) z^t
&=\frac{1}{z}\sum_{t=0}^{\infty} Q_{n_0}(n,t+1) z^{t+1}
=\frac{1}{z}\sum_{s=1}^{\infty} Q_{n_0}(n,s) z^{s}\\
&=\frac{1}{z}\left(\sum_{s=0}^{\infty} Q_{n_0}(n,s) z^{s}-Q_{n_0}(n,0)\right)
=\frac{1}{z}\left(\tilde Q_{n_0}(n,z)-\delta_{n,n_0}\right),
\end{align*}
where we used $Q_{n_0}(n,0)=\delta_{n,n_0}$.
Substituting these into the summed master equation yields
\begin{align*}
\frac{1}{z}\left(\tilde Q_{n_0}(n,z)-\delta_{n,n_0}\right)
&=\frac q2\,\tilde Q_{n_0}(n-1,z)+\frac q2\,\tilde Q_{n_0}(n+1,z)\\
&\quad +(1-q)\,\tilde Q_{n_0}(n,z),\\
\tilde Q_{n_0}(n,z)-\delta_{n,n_0}
&=\frac{qz}{2}\Bigl(\tilde Q_{n_0}(n-1,z)+\tilde Q_{n_0}(n+1,z)\Bigr)+z(1-q)\,\tilde Q_{n_0}(n,z),\\
\bigl[1-z(1-q)\bigr]\tilde Q_{n_0}(n,z)
&-\frac{qz}{2}\Bigl(\tilde Q_{n_0}(n-1,z)+\tilde Q_{n_0}(n+1,z)\Bigr)\\
&=\delta_{n,n_0}.
\end{align*}
Thus we obtain the Helmholtz-type equation
\begin{equation}
\bigl[1-z(1-q)\bigr]\tilde Q_{n_0}(n,z)-\frac{qz}{2}\Bigl(\tilde Q_{n_0}(n-1,z)+\tilde Q_{n_0}(n+1,z)\Bigr)
=\delta_{n,n_0}.
\label{eq:helmholtz_ring}
\end{equation}

Diagonalize on the ring by discrete Fourier modes. Let
\[
\theta_k\equiv \frac{2\pi k}{N},\qquad k=0,1,\dots,N-1.
\]
The eigenvalues of $\underline{B}$ are
\[
\lambda_k=1-q+q\cos\theta_k.
\]
Hence the Green's function (generating function of the propagator) is
\begin{equation}
\tilde Q_{n_0}(n,z)
=\frac{1}{N}\sum_{k=0}^{N-1}\frac{e^{i\theta_k(n-n_0)}}{1-z\lambda_k}
=\frac{1}{N}\sum_{k=0}^{N-1}\frac{e^{i\theta_k(n-n_0)}}{1-z\,[1-q+q\cos\theta_k]}.
\label{eq:Q_fourier}
\end{equation}
By translational invariance, $\tilde Q_{n_0}(n,z)$ depends only on $d=d(n,n_0)$.

\subsection{Preparing the sum for the Chebyshev identity}

Introduce the spectral parameter
\begin{equation}
\sigma(z)\equiv \frac{1-z(1-q)}{qz}.
\label{eq:sigma_def}
\end{equation}
Then
\[
1-z[1-q+q\cos\theta_k]=qz\bigl(\sigma(z)-\cos\theta_k\bigr).
\]
Insert this into \eqref{eq:Q_fourier} and use symmetry of the ring to write the numerator as a cosine:
\begin{equation}
\tilde Q_{n_0}(n,z)
=\frac{1}{qz}\,\frac{1}{N}\sum_{k=0}^{N-1}\frac{\cos\!\bigl(d\,\theta_k\bigr)}{\sigma(z)-\cos\theta_k},
\qquad d=d(n,n_0).
\label{eq:Q_cos_sum}
\end{equation}

\subsection{Applying Identity E3 (Appendix E, 2020 paper) with the $k=0$ term shown explicitly}

Identity E3 for the periodic case reads (for $m\in\{0,1,\dots,N\}$)
\begin{equation}
\sum_{k=1}^{N-1}\frac{\cos\!\left(\frac{2\pi m k}{N}\right)}{\sigma-\cos\!\left(\frac{2\pi k}{N}\right)}
=
\frac{1}{1-\sigma}
+
N\,\frac{\T_{N-m}(\sigma)+\T_{m}(\sigma)}{(\sigma^2-1)\,\U_{N-1}(\sigma)}.
\label{eq:E3_identity}
\end{equation}

We now apply \eqref{eq:E3_identity} to the sum in \eqref{eq:Q_cos_sum}.
First, we use $\theta_k=\frac{2\pi k}{N}$ to \emph{align} the cosine:
\[
\cos\!\bigl(d\,\theta_k\bigr)
=\cos\!\left(d\frac{2\pi k}{N}\right)
=\cos\!\left(\frac{2\pi (m=d)\,k}{N}\right).
\]
Thus we set
\[
m=d,\qquad \sigma=\sigma(z)
\]
in \eqref{eq:E3_identity}, which gives an expression for the \emph{partial} sum $k=1,\dots,N-1$:
\begin{equation}
\sum_{k=1}^{N-1}\frac{\cos(d\theta_k)}{\sigma(z)-\cos\theta_k}
=
\frac{1}{1-\sigma(z)}
+
N\,\frac{\T_{N-d}(\sigma(z))+\T_{d}(\sigma(z))}{(\sigma(z)^2-1)\,\U_{N-1}(\sigma(z))}.
\label{eq:partial_sum}
\end{equation}

Next, we restore the missing $k=0$ term in the full sum of \eqref{eq:Q_cos_sum}:
\[
\sum_{k=0}^{N-1}\frac{\cos(d\theta_k)}{\sigma(z)-\cos\theta_k}
=
\underbrace{\frac{\cos(0)}{\sigma(z)-\cos(0)}}_{\text{$k=0$ term}}
+
\sum_{k=1}^{N-1}\frac{\cos(d\theta_k)}{\sigma(z)-\cos\theta_k}.
\]
Since $\cos(0)=1$, the $k=0$ term is
\[
\frac{\cos(0)}{\sigma(z)-\cos(0)}=\frac{1}{\sigma(z)-1}.
\]
Using \eqref{eq:partial_sum}, we obtain
\begin{align}
\sum_{k=0}^{N-1}\frac{\cos(d\theta_k)}{\sigma(z)-\cos\theta_k}
&=
\frac{1}{\sigma(z)-1}
+
\frac{1}{1-\sigma(z)}
+
N\,\frac{\T_{N-d}(\sigma(z))+\T_{d}(\sigma(z))}{(\sigma(z)^2-1)\,\U_{N-1}(\sigma(z))}.
\label{eq:full_sum_with_cancel}
\end{align}
The first two terms cancel \emph{exactly}:
\[
\frac{1}{\sigma(z)-1}+\frac{1}{1-\sigma(z)}
=
\frac{1}{\sigma(z)-1}-\frac{1}{\sigma(z)-1}=0.
\]
Therefore the full sum becomes
\begin{equation}
\sum_{k=0}^{N-1}\frac{\cos(d\theta_k)}{\sigma(z)-\cos\theta_k}
=
N\,\frac{\T_{N-d}(\sigma(z))+\T_{d}(\sigma(z))}{(\sigma(z)^2-1)\,\U_{N-1}(\sigma(z))}.
\label{eq:full_sum_final}
\end{equation}

Finally, substitute \eqref{eq:full_sum_final} back into \eqref{eq:Q_cos_sum} to obtain
\begin{equation}
\boxed{
\tilde Q_{n_0}(n,z)
=
\frac{\T_{d}(\sigma(z))+\T_{N-d}(\sigma(z))}
{qz\,\bigl(\sigma(z)^2-1\bigr)\,\U_{N-1}(\sigma(z))},
\qquad d=d(n,n_0).
}
\label{eq:Q_chebyshev}
\end{equation}

\subsection{Odd vs even $N$ (explicit)}

Let $d=d(n,n_0)\in\{0,1,\dots,\lfloor N/2\rfloor\}$.

\paragraph{$N$ odd: $N=2M+1$.}
Then $d\in\{0,1,\dots,M\}$ and $N-d\neq d$ unless $d=0$. Equation \eqref{eq:Q_chebyshev} holds as written.

\paragraph{$N$ even: $N=2M$.}
Then $d\in\{0,1,\dots,M\}$.
For $d<M$, \eqref{eq:Q_chebyshev} holds as written.
For the antipodal case $d=M=N/2$, the two Chebyshev terms coincide. Indeed, the numerator in \eqref{eq:Q_chebyshev} is
$\T_d(\sigma(z))+\T_{N-d}(\sigma(z))$, and for even $N$ with $d=N/2$ we have $N-d=d$, hence
\[
N-d=d
\quad\Rightarrow\quad
\T_d(\sigma(z))+\T_{N-d}(\sigma(z))
=\T_{N/2}(\sigma(z))+\T_{N/2}(\sigma(z))
=2\,\T_{N/2}(\sigma(z)).
\]
Therefore,
\[
\tilde Q_{n_0}(n,z)\big|_{d=N/2}
=
\frac{2\,\T_{N/2}(\sigma(z))}
{qz\,\bigl(\sigma(z)^2-1\bigr)\,\U_{N-1}(\sigma(z))}.
\]

\section{Part 2 --- Heterogeneous propagator $\tilde S_{n_0}(n,z)$ (one directed shortcut)}

\subsection{Defect-technique formula (matrix form)}

For $M$ (inert, probability-preserving) defect pairs $(u_i,v_i)$ with defect parameters $\eta_{u_i,v_i}$ and $\eta_{v_i,u_i}$,
the generating function of the heterogeneous propagator is
\begin{equation}
\tilde S_{n_0}(n,z)
=
\tilde Q_{n_0}(n,z)-1+\frac{\det\!\bigl[\underline{H}(n,n_0;z)\bigr]}{\det\!\bigl[\underline{H}(z)\bigr]}.
\label{eq:defect_formula}
\end{equation}
The $M\times M$ matrices $\underline{H}(z)$ and $\underline{H}(n,n_0;z)$ are defined componentwise by
\begin{align}
H(z)_{i,j}
&=
\eta_{u_i,v_i}\,\tilde Q_{(u_j-v_j)}(u_i,z)
-\eta_{v_i,u_i}\,\tilde Q_{(u_j-v_j)}(v_i,z)
-\frac{\delta_{i,j}}{z},
\label{eq:H_def}\\[2mm]
H(n,n_0;z)_{i,j}
&=
H(z)_{i,j}
-\tilde Q_{(u_j-v_j)}(n,z)\Bigl[
\eta_{u_i,v_i}\,\tilde Q_{n_0}(u_i,z)
-\eta_{v_i,u_i}\,\tilde Q_{n_0}(v_i,z)
\Bigr],
\label{eq:Hn_def}
\end{align}
with the difference notation
\[
\tilde Q_{(u_j-v_j)}(x,z)\equiv \tilde Q_{u_j}(x,z)-\tilde Q_{v_j}(x,z).
\]

\subsection{Single directed shortcut (nearest-neighbour $k=1$; $K=2$ defect sites; $M=1$ defect pair)}

A single directed shortcut involves exactly two special sites, $u$ and $v$ (hence $K=2$ defect sites; this $K$ counts only the endpoints),
but it corresponds to one defect \emph{pair} in the formalism: $M=1$ with $(u_1,v_1)=(u,v)$.

For a directed shortcut $u\to v$ of probability $p$ taken from the lazy weight at $u$, the only nonzero off-diagonal defect
parameter is
\[
\eta_{u,v}=X_{v,u}=B_{v,u}-A_{v,u}=-p,\qquad \eta_{v,u}=0.
\]
Since $M=1$, determinants reduce to scalars: $\det[\underline{H}]=H$.

\subsection{Explicit scalar determinants}

Using \eqref{eq:H_def}--\eqref{eq:Hn_def} with $M=1$ gives
\begin{align}
\det[\underline{H}(z)]
&=
H(z)
=
\eta_{u,v}\,\tilde Q_{(u-v)}(u,z)-\eta_{v,u}\,\tilde Q_{(u-v)}(v,z)-\frac{1}{z},
\label{eq:H_scalar}\\[1mm]
\det[\underline{H}(n,n_0;z)]
&=
H(n,n_0;z)
=
H(z)-\tilde Q_{(u-v)}(n,z)\Bigl[\eta_{u,v}\,\tilde Q_{n_0}(u,z)-\eta_{v,u}\,\tilde Q_{n_0}(v,z)\Bigr].
\label{eq:Hn_scalar}
\end{align}

On the homogeneous ring, $\tilde Q_{a}(b,z)$ depends only on $d(a,b)$. In particular, letting $d_{uv}=d(u,v)$:
\[
\tilde Q_{(u-v)}(u,z)=\tilde Q_u(u,z)-\tilde Q_v(u,z)=\tilde Q(0,z)-\tilde Q(d_{uv},z),
\]
and
\[
\tilde Q_{(u-v)}(n,z)=\tilde Q_u(n,z)-\tilde Q_v(n,z)=\tilde Q(d(n,u),z)-\tilde Q(d(n,v),z).
\]
Also $\tilde Q_{n_0}(u,z)=\tilde Q(d(u,n_0),z)$.

With $\eta_{u,v}=-p$ and $\eta_{v,u}=0$, \eqref{eq:H_scalar}--\eqref{eq:Hn_scalar} become
\begin{align}
\boxed{
\det[\underline{H}(z)]
=
-\frac{1}{z}
-p\Bigl(\tilde Q(0,z)-\tilde Q(d_{uv},z)\Bigr),
}
\label{eq:H_final}\\[2mm]
\boxed{
\det[\underline{H}(n,n_0;z)]
=
\det[\underline{H}(z)]
+p\Bigl(\tilde Q(d(n,u),z)-\tilde Q(d(n,v),z)\Bigr)\,\tilde Q(d(u,n_0),z).
}
\label{eq:Hn_final}
\end{align}

\subsection{Closed form for $\tilde S_{n_0}(n,z)$}

Insert \eqref{eq:H_final}--\eqref{eq:Hn_final} into \eqref{eq:defect_formula}:
\begin{align*}
\tilde S_{n_0}(n,z)
&=
\tilde Q_{n_0}(n,z)-1+\frac{\det[\underline{H}(n,n_0;z)]}{\det[\underline{H}(z)]}\\
&=
\tilde Q_{n_0}(n,z)
 +
\frac{
p\Bigl(\tilde Q(d(n,u),z)-\tilde Q(d(n,v),z)\Bigr)\,\tilde Q(d(u,n_0),z)
}{
\det[\underline{H}(z)]
}.
\end{align*}
Since $\det[\underline{H}(z)]=-(1/z)\bigl[1+zp(\tilde Q(0,z)-\tilde Q(d_{uv},z))\bigr]$, we obtain the explicit rational form
\begin{equation}
\boxed{
\begin{aligned}
\tilde S_{n_0}(n,z)
&=
\tilde Q(d(n,n_0),z)\\
&\quad
- \frac{z\,p}{1+z\,p\Bigl(\tilde Q(0,z)-\tilde Q(d(u,v),z)\Bigr)}\\
&\qquad\qquad \times \Bigl(\tilde Q(d(n,u),z)-\tilde Q(d(n,v),z)\Bigr)\,\tilde Q(d(u,n_0),z).
\end{aligned}
}
\label{eq:S_closed}
\end{equation}

\subsection{Determinants in Chebyshev form}

Define, for compactness,
\[
\sigma=\sigma(z),\qquad
\mathcal{C}(z)\equiv qz\bigl(\sigma^2-1\bigr)\U_{N-1}(\sigma),\qquad
\mathcal{N}_d(\sigma)\equiv \T_d(\sigma)+\T_{N-d}(\sigma).
\]
Then \eqref{eq:Q_chebyshev} reads
\[
\tilde Q(d,z)=\frac{\mathcal{N}_d(\sigma)}{\mathcal{C}(z)}.
\]
Let $d_{ab}=d(a,b)$. Using \eqref{eq:H_final}--\eqref{eq:Hn_final} we obtain
\begin{equation}
\boxed{
\begin{aligned}
\det[\underline{H}(z)]
&=
-\frac{1}{z}
-\frac{p}{\mathcal{C}(z)}\Bigl(\mathcal{N}_0(\sigma)-\mathcal{N}_{d_{uv}}(\sigma)\Bigr)\\
&=
-\frac{1}{z}
-\frac{p}{\mathcal{C}(z)}\Bigl(1+\T_N(\sigma)-\T_{d_{uv}}(\sigma)-\T_{N-d_{uv}}(\sigma)\Bigr).
\end{aligned}
}
\label{eq:H_cheb}
\end{equation}
and
\begin{equation}
\boxed{
\det[\underline{H}(n,n_0;z)]
=
\det[\underline{H}(z)]
+\frac{p}{\mathcal{C}(z)^2}
\Bigl(\mathcal{N}_{d_{nu}}(\sigma)-\mathcal{N}_{d_{nv}}(\sigma)\Bigr)\,
\mathcal{N}_{d_{un_0}}(\sigma).
}
\label{eq:Hn_cheb}
\end{equation}

\section{Part 3 --- First-passage time (FPT) generating function}

The FPT generating function from $n_0$ to $n$ is
\begin{equation}
\tilde F_{n_0\to n}(z)=\frac{\tilde S_{n_0}(n,z)}{\tilde S_{n}(n,z)}.
\label{eq:F_def}
\end{equation}
Using \eqref{eq:S_closed} both in the numerator and with $n_0=n$ in the denominator yields
\begin{equation}
\boxed{
\begin{aligned}
\tilde F_{n_0\to n}(z)
&=
\frac{
\tilde Q(d(n,n_0),z)
-
\dfrac{
z\,p\Bigl(\tilde Q(d(n,u),z)-\tilde Q(d(n,v),z)\Bigr)\,\tilde Q(d(u,n_0),z)
}{
1+z\,p\Bigl(\tilde Q(0,z)-\tilde Q(d(u,v),z)\Bigr)
}
}{
\tilde Q(0,z)
-
\dfrac{
z\,p\Bigl(\tilde Q(d(n,u),z)-\tilde Q(d(n,v),z)\Bigr)\,\tilde Q(d(u,n),z)
}{
1+z\,p\Bigl(\tilde Q(0,z)-\tilde Q(d(u,v),z)\Bigr)
}
}.
\end{aligned}
}
\label{eq:F_final}
\end{equation}

\subsection{Algebraic simplification and Chebyshev closed form}

Starting from the nested-fraction form \eqref{eq:F_final}, define the two recurring algebraic blocks
\begin{equation}
D_Q(z)\equiv 1+z p\bigl(\tilde Q(0,z)-\tilde Q(d(u,v),z)\bigr),
\qquad
\Delta_n(z)\equiv \tilde Q(d(n,u),z)-\tilde Q(d(n,v),z).
\label{eq:DQ_Delta_def}
\end{equation}

\paragraph{Step 1: remove nested fractions.}
Multiply the numerator and denominator of \eqref{eq:F_final} by $D_Q(z)$ to obtain the equivalent form without nested fractions:
\begin{equation}
\boxed{
\tilde F_{n_0\to n}(z)=
\frac{
\tilde Q(d(n,n_0),z)\,D_Q(z)-z p\,\Delta_n(z)\,\tilde Q(d(u,n_0),z)
}{
\tilde Q(0,z)\,D_Q(z)-z p\,\Delta_n(z)\,\tilde Q(d(u,n),z)
}.
}
\label{eq:F_no_nested}
\end{equation}

\paragraph{Step 2: substitute the Chebyshev packaging for $\tilde Q(d,z)$.}
Using $\sigma(z)$ from \eqref{eq:sigma_def} and the Chebyshev form \eqref{eq:Q_chebyshev}, define
\[
\tilde Q(d,z)=\frac{\mathcal{N}_d(\sigma(z))}{\mathcal{C}(z)},
\qquad
\mathcal{C}(z)\equiv qz\bigl(\sigma(z)^2-1\bigr)\U_{N-1}(\sigma(z)),
\qquad
\mathcal{N}_d(\sigma)\equiv \T_d(\sigma)+\T_{N-d}(\sigma).
\]
Then
\[
\Delta_n(z)=\frac{\mathcal{N}_{d(n,u)}(\sigma(z))-\mathcal{N}_{d(n,v)}(\sigma(z))}{\mathcal{C}(z)},
\qquad
D_Q(z)=\frac{\mathcal{C}(z)+z p\bigl(\mathcal{N}_0(\sigma(z))-\mathcal{N}_{d(u,v)}(\sigma(z))\bigr)}{\mathcal{C}(z)}.
\]
It is convenient to introduce the $u$--$v$ dependent factor
\begin{equation}
\boxed{
D(z)\equiv \mathcal{C}(z)+z p\bigl(\mathcal{N}_0(\sigma(z))-\mathcal{N}_{d(u,v)}(\sigma(z))\bigr),
}
\label{eq:D_def}
\end{equation}
so that $D_Q(z)=D(z)/\mathcal{C}(z)$.

\paragraph{Step 3: cancel common factors and obtain the closed form.}
Substituting the above expressions into \eqref{eq:F_no_nested} shows that both the numerator and denominator carry the same factor $\mathcal{C}(z)^{-2}$, which cancels.
The result is the fully explicit closed form
\begin{equation}
\boxed{
\tilde F_{n_0\to n}(z)=
\frac{
\mathcal{N}_{d(n,n_0)}(\sigma(z))\,D(z)
-z p\,\Bigl(\mathcal{N}_{d(n,u)}(\sigma(z))-\mathcal{N}_{d(n,v)}(\sigma(z))\Bigr)\,\mathcal{N}_{d(u,n_0)}(\sigma(z))
}{
\mathcal{N}_{0}(\sigma(z))\,D(z)
-z p\,\Bigl(\mathcal{N}_{d(n,u)}(\sigma(z))-\mathcal{N}_{d(n,v)}(\sigma(z))\Bigr)\,\mathcal{N}_{d(u,n)}(\sigma(z))
}.
}
\label{eq:F_cheb_closed}
\end{equation}
Here $\mathcal{N}_0(\sigma)=\T_0(\sigma)+\T_N(\sigma)=1+\T_N(\sigma)$.

\paragraph{Remark (antipodal distance for even $N$).}
If $N$ is even and some distance equals $d=N/2$, then $\mathcal{N}_{N/2}(\sigma)=2\,\T_{N/2}(\sigma)$, and the above formulas simplify accordingly.

\section*{Numerical validation (double precision)}

We validated \eqref{eq:Q_chebyshev}, \eqref{eq:S_closed}, and \eqref{eq:F_cheb_closed} against direct matrix computation of the resolvent.
For each parameter set $(N,q,p,u,v,z)$ we built the column-stochastic transition matrices $\underline{B}$ (ring) and $\underline{A}$ (directed shortcut),
and computed
\[
G_B(z)=(\underline{I}-z\underline{B})^{-1},\qquad
G_A(z)=(\underline{I}-z\underline{A})^{-1}.
\]
The ground-truth generating functions are $\tilde Q_{n_0}(n,z)=[G_B(z)]_{n,n_0}$ and $\tilde S_{n_0}(n,z)=[G_A(z)]_{n,n_0}$.
For first passage we used $\tilde F_{n_0\to n}(z)=[G_A(z)]_{n,n_0}/[G_A(z)]_{n,n}$ for $n_0\neq n$.
We report the maximum absolute error over all indices:
\[
\max_{n,n_0}\bigl|\Delta\tilde Q\bigr|,\qquad
\max_{n,n_0}\bigl|\Delta\tilde S\bigr|,\qquad
\max_{n_0\neq n}\bigl|\Delta\tilde F\bigr|.
\]

\begin{center}
\small
\begin{tabular}{r r r r r r l l l}
\hline
$N$ & $q$ & $p$ & $u$ & $v$ & $z$ & $\max|\Delta\tilde Q|$ & $\max|\Delta\tilde S|$ & $\max|\Delta\tilde F|$ \\
\hline
9 & 0.6 & 0.2 & 1 & 4 & 0.3 & 2.22e-16 & 4.44e-16 & 5.55e-17 \\
10 & 0.7 & 0.1 & 2 & 7 & 0.5 & 4.44e-16 & 2.22e-16 & 5.55e-17 \\
12 & 0.5 & 0.3 & 0 & 6 & 0.6 & 4.44e-16 & 4.44e-16 & 8.33e-17 \\
15 & 0.8 & 0.05 & 3 & 4 & 0.4 & 2.22e-16 & 4.44e-16 & 1.11e-16 \\
20 & 0.4 & 0.2 & 5 & 13 & 0.75 & 1.33e-15 & 1.33e-15 & 1.67e-16 \\
\hline
\end{tabular}
\end{center}

In a sweep of 200 random cases ($N\le 40$, seed $=0$, $q\in[0.2,0.9]$, $z\in[0.2,0.85]$, $p\in[0,1-q]$), we found
\[
\max|\Delta\tilde Q|=3.11\times 10^{-15},\qquad
\max|\Delta\tilde S|=3.11\times 10^{-15},\qquad
\max|\Delta\tilde F|=4.44\times 10^{-16}.
\]
These residuals are consistent with floating-point round-off and the conditioning of matrix inversion as $z\to 1^{-}$.

\end{document}
